% TOP_PROBLEM: {"id":30007570,"problem_number":"LOCAL-30007570","title":"Statistical confidence in game equilibria","category_id":21,"category":{"id":21,"name":"economics","display_name":"Economics","description":"Open economic theory statements, published research agendas, and proposed scoped specifications; question provenance and historical priority are qualified per record.","slug":"economics","order_index":21,"created_at":"2026-10-08T00:00:00Z"},"status":"research_question","external_url":null,"legacy_ids":[],"published":false,"statement":"Design cost-efficient adaptive simulation and strategy exploration that outputs a statistically valid full-game approximate-equilibrium certificate despite optional stopping and candidate reuse.","economics":{"original_id":59,"field":"Game theory and strategic interaction","kind":"design","formalization_basis":"adapted_research_question","source_status":"research_agenda","priority_status":"earliest_found","priority_note":"This is the earliest source in which the relevant agenda was verified during this audit; absolute priority has not been established.","earliest_verified_reference":{"authors":["Michael P. Wellman","Karl Tuyls","Amy Greenwald"],"title":"Empirical Game-Theoretic Analysis: A Survey","year":2025,"url":"https://arxiv.org/pdf/2403.04018","doi":null,"locator":"§6 Statistical Reasoning in Empirical Games, pp.1045–1052; §7.3 Frontiers in Strategy Exploration, pp.1057–1058","evidence_summary":"The survey discusses adaptive statistical conclusions and identifies unresolved theoretical characterization of strategy-exploration solvers."},"current_reference":{"authors":["Michael P. Wellman","Karl Tuyls","Amy Greenwald"],"title":"Empirical Game-Theoretic Analysis: A Survey","year":2025,"url":"https://arxiv.org/pdf/2403.04018","doi":null,"locator":"§6 Statistical Reasoning in Empirical Games, pp.1045–1052; §7.3 Frontiers in Strategy Exploration, pp.1057–1058","evidence_summary":"The survey discusses adaptive statistical conclusions and identifies unresolved theoretical characterization of strategy-exploration solvers."},"formalization_note":"The joint adaptive certificate and cost target are introduced here. Existing statistical EGTA methods must be incorporated when determining residual novelty."},"statusQualification":"Published question or research agenda verified at the cited locator. The mathematical specification is adapted for this catalogue and includes additional assumptions; source publication does not establish first-ever priority or certify this exact model remains unsolved. The joint adaptive certificate and cost target are introduced here. Existing statistical EGTA methods must be incorporated when determining residual novelty.","statusReviewedAt":"2026-10-08"}
% FIRST_POSTED: 2026-10-08
% ENTRY_KIND: statement_only
% !TeX program = lualatex
\documentclass[11pt]{article}
\usepackage{fontspec}
\usepackage[a4paper,margin=25mm]{geometry}
\usepackage{amsmath,amssymb,mathtools}
\usepackage{xurl}
\usepackage[hidelinks]{hyperref}
\newcommand{\sref}[2]{\hyperlink{source:#1:#2}{[S#2]}}
\setlength{\emergencystretch}{3em}
\begin{document}
% problemId:30007570
\section*{Statistical confidence in game equilibria}\label{30007570}
\subsection{Definitions and mathematical statement}
Let \(G\) be a finite \(n\)-player game with action sets \(A_i\) and unknown expected payoffs \(u_i(a)\in[0,1]\). A simulator queried at profile \(a\) returns independent bounded payoff samples with known query cost \(c(a)>0\); samples may have profile-dependent variance. An adaptive procedure chooses simulator queries, constructs candidate mixed profiles from the same observations, and stops at a random time with output \((\widehat\sigma,\widehat e)\). Define full-game exploitability
\[
 e_G(\sigma)=\max_i\max_{b_i\in A_i}
 \{u_i(b_i,\sigma_{-i})-u_i(\sigma)\}.
\]
For tolerances \(\varepsilon>0\) and \(\delta\in(0,1)\), require a simultaneously valid certificate \(\Pr(e_G(\widehat\sigma)\leq\widehat e\leq\varepsilon)\geq1-\delta\), uniformly over the maintained payoff class. Determine sharp expected simulation-cost bounds and adaptive sampling rules that approach them, accounting for candidate selection, optional stopping, and unevaluated deviations. State how bounds depend on action count, variance, payoff gaps, and exploitable structure. Exhaustively estimating every payoff is an available baseline; the research target is principled savings while preserving certification as strategy exploration changes the candidate. A valid solution must also specify failure or nontermination reporting when a prescribed computation budget is insufficient. This is a scoped adaptation of the survey’s statistical and exploration agendas, not a claim that confidence bounds for fixed finite games are unknown.

\par\textbf{Formulation and scope.} The joint adaptive certificate and cost target are introduced here. Existing statistical EGTA methods must be incorporated when determining residual novelty.

\par\textbf{Open-question provenance.} An explicit question or research agenda was verified at the source locator. This does not by itself certify that every later version remains unresolved \sref{30007570}{1}.

\par\textbf{Historical priority.} This is the earliest source in which the relevant agenda was verified during this audit; absolute priority has not been established.

\subsection{Short English statement}
Design cost-efficient adaptive simulation and strategy exploration that outputs a statistically valid full-game approximate-equilibrium certificate despite optional stopping and candidate reuse.

\subsection{Sources}
\begin{itemize}\raggedright
\item[\textnormal{[S1]}]\hypertarget{source:30007570:1}{} Michael P. Wellman, Karl Tuyls, Amy Greenwald. 2025. Empirical Game-Theoretic Analysis: A Survey. §6 Statistical Reasoning in Empirical Games, pp.1045–1052; §7.3 Frontiers in Strategy Exploration, pp.1057–1058.
\url{https://arxiv.org/pdf/2403.04018}.
{\small\textit{Earliest verified question/agenda reference; historical priority qualified below. The survey discusses adaptive statistical conclusions and identifies unresolved theoretical characterization of strategy-exploration solvers.}}
\end{itemize}
\end{document}
